\IfFileExists{stacks-project.cls}{%
\documentclass{stacks-project}
}{%
\documentclass{amsart}
}

% For dealing with references we use the comment environment
\usepackage{verbatim}
\newenvironment{reference}{\comment}{\endcomment}
%\newenvironment{reference}{}{}
\newenvironment{slogan}{\comment}{\endcomment}
\newenvironment{history}{\comment}{\endcomment}

% For commutative diagrams we use Xy-pic
\usepackage[all]{xy}

% We use 2cell for 2-commutative diagrams.
\xyoption{2cell}
\UseAllTwocells

% We use multicol for the list of chapters between chapters
\usepackage{multicol}

% This is generally recommended for better output
\usepackage{lmodern}
\usepackage[T1]{fontenc}

% For cross-file-references
\usepackage{xr-hyper}

% Package for hypertext links:
\usepackage{hyperref}

% For any local file, say "hello.tex" you want to link to please
% use \externaldocument[hello-]{hello}
\externaldocument[introduction-]{introduction}
\externaldocument[conventions-]{conventions}
\externaldocument[sets-]{sets}
\externaldocument[categories-]{categories}
\externaldocument[topology-]{topology}
\externaldocument[sheaves-]{sheaves}
\externaldocument[sites-]{sites}
\externaldocument[stacks-]{stacks}
\externaldocument[fields-]{fields}
\externaldocument[algebra-]{algebra}
\externaldocument[brauer-]{brauer}
\externaldocument[homology-]{homology}
\externaldocument[derived-]{derived}
\externaldocument[simplicial-]{simplicial}
\externaldocument[more-algebra-]{more-algebra}
\externaldocument[smoothing-]{smoothing}
\externaldocument[modules-]{modules}
\externaldocument[sites-modules-]{sites-modules}
\externaldocument[injectives-]{injectives}
\externaldocument[cohomology-]{cohomology}
\externaldocument[sites-cohomology-]{sites-cohomology}
\externaldocument[dga-]{dga}
\externaldocument[dpa-]{dpa}
\externaldocument[sdga-]{sdga}
\externaldocument[hypercovering-]{hypercovering}
\externaldocument[schemes-]{schemes}
\externaldocument[constructions-]{constructions}
\externaldocument[properties-]{properties}
\externaldocument[morphisms-]{morphisms}
\externaldocument[coherent-]{coherent}
\externaldocument[divisors-]{divisors}
\externaldocument[limits-]{limits}
\externaldocument[varieties-]{varieties}
\externaldocument[topologies-]{topologies}
\externaldocument[descent-]{descent}
\externaldocument[perfect-]{perfect}
\externaldocument[more-morphisms-]{more-morphisms}
\externaldocument[flat-]{flat}
\externaldocument[groupoids-]{groupoids}
\externaldocument[more-groupoids-]{more-groupoids}
\externaldocument[etale-]{etale}
\externaldocument[chow-]{chow}
\externaldocument[intersection-]{intersection}
\externaldocument[pic-]{pic}
\externaldocument[weil-]{weil}
\externaldocument[adequate-]{adequate}
\externaldocument[dualizing-]{dualizing}
\externaldocument[duality-]{duality}
\externaldocument[discriminant-]{discriminant}
\externaldocument[derham-]{derham}
\externaldocument[local-cohomology-]{local-cohomology}
\externaldocument[algebraization-]{algebraization}
\externaldocument[curves-]{curves}
\externaldocument[resolve-]{resolve}
\externaldocument[models-]{models}
\externaldocument[functors-]{functors}
\externaldocument[equiv-]{equiv}
\externaldocument[pione-]{pione}
\externaldocument[etale-cohomology-]{etale-cohomology}
\externaldocument[proetale-]{proetale}
\externaldocument[relative-cycles-]{relative-cycles}
\externaldocument[more-etale-]{more-etale}
\externaldocument[trace-]{trace}
\externaldocument[crystalline-]{crystalline}
\externaldocument[spaces-]{spaces}
\externaldocument[spaces-properties-]{spaces-properties}
\externaldocument[spaces-morphisms-]{spaces-morphisms}
\externaldocument[decent-spaces-]{decent-spaces}
\externaldocument[spaces-cohomology-]{spaces-cohomology}
\externaldocument[spaces-limits-]{spaces-limits}
\externaldocument[spaces-divisors-]{spaces-divisors}
\externaldocument[spaces-over-fields-]{spaces-over-fields}
\externaldocument[spaces-topologies-]{spaces-topologies}
\externaldocument[spaces-descent-]{spaces-descent}
\externaldocument[spaces-perfect-]{spaces-perfect}
\externaldocument[spaces-more-morphisms-]{spaces-more-morphisms}
\externaldocument[spaces-flat-]{spaces-flat}
\externaldocument[spaces-groupoids-]{spaces-groupoids}
\externaldocument[spaces-more-groupoids-]{spaces-more-groupoids}
\externaldocument[bootstrap-]{bootstrap}
\externaldocument[spaces-pushouts-]{spaces-pushouts}
\externaldocument[spaces-chow-]{spaces-chow}
\externaldocument[groupoids-quotients-]{groupoids-quotients}
\externaldocument[spaces-more-cohomology-]{spaces-more-cohomology}
\externaldocument[spaces-simplicial-]{spaces-simplicial}
\externaldocument[spaces-duality-]{spaces-duality}
\externaldocument[formal-spaces-]{formal-spaces}
\externaldocument[restricted-]{restricted}
\externaldocument[spaces-resolve-]{spaces-resolve}
\externaldocument[formal-defos-]{formal-defos}
\externaldocument[defos-]{defos}
\externaldocument[cotangent-]{cotangent}
\externaldocument[examples-defos-]{examples-defos}
\externaldocument[algebraic-]{algebraic}
\externaldocument[examples-stacks-]{examples-stacks}
\externaldocument[stacks-sheaves-]{stacks-sheaves}
\externaldocument[criteria-]{criteria}
\externaldocument[artin-]{artin}
\externaldocument[quot-]{quot}
\externaldocument[stacks-properties-]{stacks-properties}
\externaldocument[stacks-morphisms-]{stacks-morphisms}
\externaldocument[stacks-limits-]{stacks-limits}
\externaldocument[stacks-cohomology-]{stacks-cohomology}
\externaldocument[stacks-perfect-]{stacks-perfect}
\externaldocument[stacks-introduction-]{stacks-introduction}
\externaldocument[stacks-more-morphisms-]{stacks-more-morphisms}
\externaldocument[stacks-geometry-]{stacks-geometry}
\externaldocument[moduli-]{moduli}
\externaldocument[moduli-curves-]{moduli-curves}
\externaldocument[examples-]{examples}
\externaldocument[exercises-]{exercises}
\externaldocument[guide-]{guide}
\externaldocument[desirables-]{desirables}
\externaldocument[coding-]{coding}
\externaldocument[obsolete-]{obsolete}
\externaldocument[fdl-]{fdl}
\externaldocument[index-]{index}

% Theorem environments.
%
\theoremstyle{plain}
\newtheorem{theorem}[subsection]{Theorem}
\newtheorem{proposition}[subsection]{Proposition}
\newtheorem{lemma}[subsection]{Lemma}

\theoremstyle{definition}
\newtheorem{definition}[subsection]{Definition}
\newtheorem{example}[subsection]{Example}
\newtheorem{exercise}[subsection]{Exercise}
\newtheorem{situation}[subsection]{Situation}

\theoremstyle{remark}
\newtheorem{remark}[subsection]{Remark}
\newtheorem{remarks}[subsection]{Remarks}

\numberwithin{equation}{subsection}

% Macros
%
\def\lim{\mathop{\mathrm{lim}}\nolimits}
\def\colim{\mathop{\mathrm{colim}}\nolimits}
\def\Spec{\mathop{\mathrm{Spec}}}
\def\Hom{\mathop{\mathrm{Hom}}\nolimits}
\def\Ext{\mathop{\mathrm{Ext}}\nolimits}
\def\SheafHom{\mathop{\mathcal{H}\!\mathit{om}}\nolimits}
\def\SheafExt{\mathop{\mathcal{E}\!\mathit{xt}}\nolimits}
\def\Sch{\mathit{Sch}}
\def\Mor{\mathop{\mathrm{Mor}}\nolimits}
\def\Ob{\mathop{\mathrm{Ob}}\nolimits}
\def\Sh{\mathop{\mathit{Sh}}\nolimits}
\def\NL{\mathop{N\!L}\nolimits}
\def\CH{\mathop{\mathrm{CH}}\nolimits}
\def\proetale{{pro\text{-}\acute{e}tale}}
\def\etale{{\acute{e}tale}}
\def\QCoh{\mathit{QCoh}}
\def\Ker{\mathop{\mathrm{Ker}}}
\def\Im{\mathop{\mathrm{Im}}}
\def\Coker{\mathop{\mathrm{Coker}}}
\def\Coim{\mathop{\mathrm{Coim}}}

% Boxtimes
%
\DeclareMathSymbol{\boxtimes}{\mathbin}{AMSa}{"02}

%
% Macros for moduli stacks/spaces
%
\def\QCohstack{\mathcal{QC}\!\mathit{oh}}
\def\Cohstack{\mathcal{C}\!\mathit{oh}}
\def\Spacesstack{\mathcal{S}\!\mathit{paces}}
\def\Quotfunctor{\mathrm{Quot}}
\def\Hilbfunctor{\mathrm{Hilb}}
\def\Curvesstack{\mathcal{C}\!\mathit{urves}}
\def\Polarizedstack{\mathcal{P}\!\mathit{olarized}}
\def\Complexesstack{\mathcal{C}\!\mathit{omplexes}}
% \Pic is the operator that assigns to X its picard group, usage \Pic(X)
% \Picardstack_{X/B} denotes the Picard stack of X over B
% \Picardfunctor_{X/B} denotes the Picard functor of X over B
\def\Pic{\mathop{\mathrm{Pic}}\nolimits}
\def\Picardstack{\mathcal{P}\!\mathit{ic}}
\def\Picardfunctor{\mathrm{Pic}}
\def\Deformationcategory{\mathcal{D}\!\mathit{ef}}
\usepackage{cleveref}
\crefname{originalsection}{original section}{original sections}
\Crefname{originalsection}{Original section}{Original sections}
\makeatletter
\newlabel{derived-section-homotopy}{{05RN}{source}{}{sourcepointer-1}{}}
\newlabel{derived-section-homotopy@cref}{{[originalsection][0][]05RN}{[1][0][]source}{}{}{}}
\makeatother
\crefname{originaldefinition}{original definition}{original definitions}
\Crefname{originaldefinition}{Original definition}{Original definitions}
\makeatletter
\newlabel{derived-definition-triangulated-subcategory}{{05QM}{source}{}{sourcepointer-2}{}}
\newlabel{derived-definition-triangulated-subcategory@cref}{{[originaldefinition][0][]05QM}{[1][0][]source}{}{}{}}
\makeatother
\crefname{originaldefinition}{original definition}{original definitions}
\Crefname{originaldefinition}{Original definition}{Original definitions}
\makeatletter
\newlabel{derived-definition-saturated}{{05RB}{source}{}{sourcepointer-3}{}}
\newlabel{derived-definition-saturated@cref}{{[originaldefinition][0][]05RB}{[1][0][]source}{}{}{}}
\makeatother
\crefname{originaldefinition}{original definition}{original definitions}
\Crefname{originaldefinition}{Original definition}{Original definitions}
\makeatletter
\newlabel{derived-definition-localization}{{05R2}{source}{}{sourcepointer-4}{}}
\newlabel{derived-definition-localization@cref}{{[originaldefinition][0][]05R2}{[1][0][]source}{}{}{}}
\makeatother
\crefname{originaldefinition}{original definition}{original definitions}
\Crefname{originaldefinition}{Original definition}{Original definitions}
\makeatletter
\newlabel{categories-definition-saturated-multiplicative-system}{{05Q8}{source}{}{sourcepointer-5}{}}
\newlabel{categories-definition-saturated-multiplicative-system@cref}{{[originaldefinition][0][]05Q8}{[1][0][]source}{}{}{}}
\makeatother
\crefname{originallemma}{original lemma}{original lemmas}
\Crefname{originallemma}{Original lemma}{Original lemmas}
\makeatletter
\newlabel{cohomology-lemma-two-out-of-three-perfect}{{08CR}{source}{}{sourcepointer-6}{}}
\newlabel{cohomology-lemma-two-out-of-three-perfect@cref}{{[originallemma][0][]08CR}{[1][0][]source}{}{}{}}
\makeatother
\crefname{originallemma}{original lemma}{original lemmas}
\Crefname{originallemma}{Original lemma}{Original lemmas}
\makeatletter
\newlabel{cohomology-lemma-summands-perfect}{{08CS}{source}{}{sourcepointer-7}{}}
\newlabel{cohomology-lemma-summands-perfect@cref}{{[originallemma][0][]08CS}{[1][0][]source}{}{}{}}
\makeatother
\crefname{originallemma}{original lemma}{original lemmas}
\Crefname{originallemma}{Original lemma}{Original lemmas}
\makeatletter
\newlabel{derived-lemma-bounded-triangulated-subcategories}{{05RQ}{source}{}{sourcepointer-8}{}}
\newlabel{derived-lemma-bounded-triangulated-subcategories@cref}{{[originallemma][0][]05RQ}{[1][0][]source}{}{}{}}
\makeatother
\crefname{originallemma}{original lemma}{original lemmas}
\Crefname{originallemma}{Original lemma}{Original lemmas}
\makeatletter
\newlabel{derived-lemma-triangle-functor-localize}{{05R4}{source}{}{sourcepointer-9}{}}
\newlabel{derived-lemma-triangle-functor-localize@cref}{{[originallemma][0][]05R4}{[1][0][]source}{}{}{}}
\makeatother
\crefname{originalproposition}{original proposition}{original propositions}
\Crefname{originalproposition}{Original proposition}{Original propositions}
\makeatletter
\newlabel{derived-proposition-construct-localization}{{05R6}{source}{}{sourcepointer-10}{}}
\newlabel{derived-proposition-construct-localization@cref}{{[originalproposition][0][]05R6}{[1][0][]source}{}{}{}}
\makeatother
\crefname{originallemma}{original lemma}{original lemmas}
\Crefname{originallemma}{Original lemma}{Original lemmas}
\makeatletter
\newlabel{derived-lemma-universal-property-localization}{{05R7}{source}{}{sourcepointer-11}{}}
\newlabel{derived-lemma-universal-property-localization@cref}{{[originallemma][0][]05R7}{[1][0][]source}{}{}{}}
\makeatother
\crefname{originallemma}{original lemma}{original lemmas}
\Crefname{originallemma}{Original lemma}{Original lemmas}
\makeatletter
\newlabel{lemma-resolution-property-perfect-complex}{{0F8F}{source}{}{sourcepointer-12}{}}
\newlabel{lemma-resolution-property-perfect-complex@cref}{{[originallemma][0][]0F8F}{[1][0][]source}{}{}{}}
\makeatother
\crefname{originallemma}{original lemma}{original lemmas}
\Crefname{originallemma}{Original lemma}{Original lemmas}
\makeatletter
\newlabel{lemma-resolution-property-map-perfect-complex}{{0F8G}{source}{}{sourcepointer-13}{}}
\newlabel{lemma-resolution-property-map-perfect-complex@cref}{{[originallemma][0][]0F8G}{[1][0][]source}{}{}{}}
\makeatother
\crefname{originallemma}{original lemma}{original lemmas}
\Crefname{originallemma}{Original lemma}{Original lemmas}
\makeatletter
\newlabel{lemma-resolution-property-homotopy-map-perfect-complex}{{0F8H}{source}{}{sourcepointer-14}{}}
\newlabel{lemma-resolution-property-homotopy-map-perfect-complex@cref}{{[originallemma][0][]0F8H}{[1][0][]source}{}{}{}}
\makeatother
\crefname{originallemma}{original lemma}{original lemmas}
\Crefname{originallemma}{Original lemma}{Original lemmas}
\makeatletter
\newlabel{derived-lemma-exact-equivalence}{{05SQ}{source}{}{sourcepointer-15}{}}
\newlabel{derived-lemma-exact-equivalence@cref}{{[originallemma][0][]05SQ}{[1][0][]source}{}{}{}}
\makeatother
\title{Perfect complexes on a scheme with the resolution property are the localization of the homotopy category of finite locally free complexes}
\author{The Stacks Project Authors}
\date{stacks-project:perfect.tex:local-snapshot-20261012; GNU FDL 1.2 or later}
\begin{document}
\maketitle
\setcounter{subsection}{1}
\noindent\textit{Editorial note.} All delivered mathematical text is the author's own, copied from the pinned Stacks Project chapter file perfect.tex; no statement or proof step was rewritten, reordered or removed. The counted span is the author environment and its own proof environment. References inside the span are handed to the assembler as navigation stubs that carry the canonical tag of the cited tag only; whether each of them is load-bearing is recorded on the stub itself.\par\medskip
\noindent\textit{Allowed standard prerequisites (editorial).}\begin{itemize}
\item The perfect complexes form a saturated triangulated subcategory of the derived category, closed under two-out-of-three and direct summands. invoked at proof lines 10023-10027. Supplies the target side of the localization statement.
\item For an exact functor out of a triangulated category, the quasi-isomorphisms in the source form a saturated multiplicative system compatible with the triangulated structure, the localization exists, and the functor factors through it universally. invoked at proof lines 10030-10042. Builds the localization S\textasciicircum{}\{-1\}K\textasciicircum{}b(A) and the exact factorization.
\item The resolution property gives: every perfect complex is represented by a bounded complex of finite locally free modules, every map between such complexes in the derived category lifts to an actual map after resolution, and equal maps become homotopic after further resolution. invoked at proof lines 10043-10052. Shows the factorized functor is an equivalence.
\end{itemize}\medskip
\ifdefined\hypertarget\hypertarget{sourcepointer-1}{}\fi
\noindent\textit{Source reference.} Original section 05RN: sources/stacks-chow/tags.txt, source lines 1--1. Referenced at statement lines 9994-9997 for the definition of K\textasciicircum{}b(A); the triangulated structure checks are inside the counted span.\par\medskip
\ifdefined\hypertarget\hypertarget{sourcepointer-2}{}\fi
\noindent\textit{Source reference.} Original definition 05QM: sources/stacks-chow/tags.txt, source lines 1--1. Cited at proof lines 10018-10021 among the definitions used; the application to D\_perf is inside the counted span.\par\medskip
\ifdefined\hypertarget\hypertarget{sourcepointer-3}{}\fi
\noindent\textit{Source reference.} Original definition 05RB: sources/stacks-chow/tags.txt, source lines 1--1. Cited at proof lines 10018-10021 among the definitions used; the verification that D\_perf is saturated is inside the counted span.\par\medskip
\ifdefined\hypertarget\hypertarget{sourcepointer-4}{}\fi
\noindent\textit{Source reference.} Original definition 05R2: sources/stacks-chow/tags.txt, source lines 1--1. Cited at proof lines 10018-10022 among the definitions used; the verification for S is inside the counted span.\par\medskip
\ifdefined\hypertarget\hypertarget{sourcepointer-5}{}\fi
\noindent\textit{Source reference.} Original definition 05Q8: sources/stacks-chow/tags.txt, source lines 1--1. Cited at proof lines 10018-10023 among the definitions used; the verification for S is inside the counted span.\par\medskip
\ifdefined\hypertarget\hypertarget{sourcepointer-6}{}\fi
\noindent\textit{Source reference.} Original lemma 08CR: sources/stacks-chow/tags.txt, source lines 1--1. Cited at proof lines 10024-10026 for the triangulated subcategory property; the localization assembly is inside the counted span.\par\medskip
\ifdefined\hypertarget\hypertarget{sourcepointer-7}{}\fi
\noindent\textit{Source reference.} Original lemma 08CS: sources/stacks-chow/tags.txt, source lines 1--1. Cited at proof lines 10024-10027 for the saturated property; the localization assembly is inside the counted span.\par\medskip
\ifdefined\hypertarget\hypertarget{sourcepointer-8}{}\fi
\noindent\textit{Source reference.} Original lemma 05RQ: sources/stacks-chow/tags.txt, source lines 1--1. Cited at proof lines 10027-10029 for K\textasciicircum{}b(A) being triangulated; the exactness of the functor is inside the counted span.\par\medskip
\ifdefined\hypertarget\hypertarget{sourcepointer-9}{}\fi
\noindent\textit{Source reference.} Original lemma 05R4: sources/stacks-chow/tags.txt, source lines 1--1. Cited at proof lines 10031-10034 for S being a compatible saturated multiplicative system; the localization construction is inside the counted span.\par\medskip
\ifdefined\hypertarget\hypertarget{sourcepointer-10}{}\fi
\noindent\textit{Source reference.} Original proposition 05R6: sources/stacks-chow/tags.txt, source lines 1--1. Cited at proof lines 10035-10038 for the existence of S\textasciicircum{}\{-1\}K\textasciicircum{}b(A); the factorization is inside the counted span.\par\medskip
\ifdefined\hypertarget\hypertarget{sourcepointer-11}{}\fi
\noindent\textit{Source reference.} Original lemma 05R7: sources/stacks-chow/tags.txt, source lines 1--1. Cited at proof lines 10039-10042 for the exact factorization; the equivalence proof is inside the counted span.\par\medskip
\ifdefined\hypertarget\hypertarget{sourcepointer-12}{}\fi
\noindent\textit{Source reference.} Original lemma 0F8F: sources/stacks-perfect/pinned-original-source.tex, source lines 9890--9890. Cited at proof lines 10043-10045 as the first equivalence input; the assembly of the three inputs is inside the counted span.\par\medskip
\ifdefined\hypertarget\hypertarget{sourcepointer-13}{}\fi
\noindent\textit{Source reference.} Original lemma 0F8G: sources/stacks-perfect/pinned-original-source.tex, source lines 9912--9912. Cited at proof lines 10043-10046 as the second equivalence input; the assembly is inside the counted span.\par\medskip
\ifdefined\hypertarget\hypertarget{sourcepointer-14}{}\fi
\noindent\textit{Source reference.} Original lemma 0F8H: sources/stacks-perfect/pinned-original-source.tex, source lines 9952--9952. Cited at proof lines 10043-10047 as the third equivalence input; the conclusion that the factorized functor is an equivalence is inside the counted span.\par\medskip
\ifdefined\hypertarget\hypertarget{sourcepointer-15}{}\fi
\noindent\textit{Source reference.} Original lemma 05SQ: sources/stacks-chow/tags.txt, source lines 1--1. Cited at proof lines 10053-10057 to conclude the equivalence is one of triangulated categories; the final assembly is inside the counted span.\par\medskip
\begin{proposition}
\label{proposition-perfect-resolution-property}
Let $X$ be a quasi-compact and quasi-separated scheme with the
resolution property. Denote
\begin{enumerate}
\item $\mathcal{A}$ the additive category of finite locally free
$\mathcal{O}_X$-modules,
\item $K^b(\mathcal{A})$ the homotopy category of bounded complexes
in $\mathcal{A}$, see
Derived Categories, Section \ref{derived-section-homotopy}, and
\item $D_{perf}(\mathcal{O}_X)$ the strictly full, saturated,
triangulated subcategory of $D(\mathcal{O}_X)$ consisting of
perfect objects.
\end{enumerate}
With this notation the obvious functor
$$
K^b(\mathcal{A}) \longrightarrow D_{perf}(\mathcal{O}_X)
$$
is an exact functor of trianglated categories which factors through an
equivalence $S^{-1}K^b(\mathcal{A}) \to D_{perf}(\mathcal{O}_X)$
of triangulated categories
where $S$ is the saturated multiplicative system of quasi-isomorphisms
in $K^b(\mathcal{A})$.
\end{proposition}

\begin{proof}
If you can parse the statement of the proposition, then please skip this
first paragraph. For some of the definitions used, please see
Derived Categories, Definition
\ref{derived-definition-triangulated-subcategory}
(triangulated subcategory),
Derived Categories, Definition \ref{derived-definition-saturated}
(saturated triangulated subcategory),
Derived Categories, Definition \ref{derived-definition-localization}
(multiplicative system compatible with the triangulated structure), and
Categories,
Definition \ref{categories-definition-saturated-multiplicative-system}
(saturated multiplicative system).
Observe that $D_{perf}(\mathcal{O}_X)$ is a saturated triangulated subcategory
of $D(\mathcal{O}_X)$ by
Cohomology, Lemmas \ref{cohomology-lemma-two-out-of-three-perfect} and
\ref{cohomology-lemma-summands-perfect}. Also, note that
$K^b(\mathcal{A})$ is a triangulated category, see
Derived Categories, Lemma
\ref{derived-lemma-bounded-triangulated-subcategories}.

\medskip\noindent
It is clear that the functor sends distinguished triangles to
distinguished triangles, i.e., is exact. Then $S$ is a saturated
multiplicative system compatible with the triangulated structure
on $K^b(\mathcal{A})$ by
Derived Categories, Lemma \ref{derived-lemma-triangle-functor-localize}.
Hence the localization $S^{-1}K^b(\mathcal{A})$ exists and is
a triangulated category by
Derived Categories, Proposition
\ref{derived-proposition-construct-localization}.
We get an exact factorization
$S^{-1}K^b(\mathcal{A}) \to D_{perf}(\mathcal{O}_X)$ by
Derived Categories, Lemma
\ref{derived-lemma-universal-property-localization}.
By Lemmas \ref{lemma-resolution-property-perfect-complex},
\ref{lemma-resolution-property-map-perfect-complex}, and
\ref{lemma-resolution-property-homotopy-map-perfect-complex}
this functor is an equivalence. Then finally the functor
$S^{-1}K^b(\mathcal{A}) \to D_{perf}(\mathcal{O}_X)$
is an equivalence of triangulated categories (in the sense that
distinguished triangles correspond) by
Derived Categories, Lemma \ref{derived-lemma-exact-equivalence}.
\end{proof}

\end{document}
